\documentclass[12pt, a4paper]{report}

\usepackage[utf8]{inputenc}
\usepackage[T1]{fontenc}
\usepackage[english]{babel}
\usepackage{amsmath, amssymb}
\usepackage{graphicx}
\usepackage{booktabs}
\usepackage{hyperref}
\usepackage{geometry}
\usepackage{setspace}
\usepackage{caption}
\usepackage{subcaption}
\usepackage{float}
\usepackage{cite}
\usepackage{url}
\usepackage{tcolorbox}  % for colored boxes in appendix

\geometry{
    a4paper,
    top=2.5cm,
    bottom=2.5cm,
    left=2.5cm,
    right=2.5cm
}

\hypersetup{
    colorlinks=false,
    linkbordercolor={0 0 0},
    citebordercolor={0 0 0}
}

\onehalfspacing

% ─────────────────────────────────────────────────────────────────────────────
\begin{document}

% ── Title page ───────────────────────────────────────────────────────────────
\begin{titlepage}
\centering
\vspace*{1cm}

{\scshape\Large Department of Mathematics\\and Computer Science\par}
\vspace{0.5cm}
{\scshape\large University of Southern Denmark\par}
{\scshape\large Bachelor Thesis in Artificial Intelligence\par}

\vspace{1.5cm}
\noindent\rule{\textwidth}{0.4pt}
{\huge\bfseries Approximate Similarity Search:\\[0.3em]
\large Evaluation on Image Retrieval, Out-of-Distribution Detection,\\
and Recommender Systems\par}
\noindent\rule{\textwidth}{0.4pt}

\vspace{2cm}

\begin{center}
\itshape Author(s)\\
\upshape Emil [Efternavn]\\
[email]@student.sdu.dk
\end{center}

\vspace{1cm}

\begin{center}
\itshape Supervisor\\
\upshape Ricardo J. G. B. Campello\\
Professor
\end{center}

\vfill
{\large May 2026\par}

\vspace{1cm}
% SDU logo — upload sdu_logo.png to Overleaf and uncomment:
% \includegraphics[width=0.25\textwidth]{sdu_logo}

\end{titlepage}

% ── Abstract ─────────────────────────────────────────────────────────────────
\begin{abstract}
\paragraph{English}
Your abstract in English.

\paragraph{Danish}
Your abstract in Danish.
\end{abstract}
\newpage

% ── Acknowledgments (optional) ───────────────────────────────────────────────
\chapter*{Acknowledgments}
Your acknowledgments, if any (optional).
\newpage

% ── Table of contents ────────────────────────────────────────────────────────
\tableofcontents
\newpage

% ─────────────────────────────────────────────────────────────────────────────
\chapter{Introduction}
% ─────────────────────────────────────────────────────────────────────────────
% Key elements: Contextualization, Motivation, Goals and Contributions.
% Suggested length: ~2 pages.



% ─────────────────────────────────────────────────────────────────────────────
\chapter{Research Question(s) and Hypothesis(es)}
% ─────────────────────────────────────────────────────────────────────────────
% What are the specific research questions?
% What hypotheses will you attempt to confirm or reject?
% Suggested length: ~1-2 pages.



% ─────────────────────────────────────────────────────────────────────────────
\chapter{Methods and Materials}
% ─────────────────────────────────────────────────────────────────────────────
% Concise summary of datasets, algorithms, experimental setup, and evaluation.
% Include detailed specifications of your free choices (esp. recommendation task).
% Brief high-level description (~0.5 page) of each ANN algorithm with hyperparameters.
% Suggested length: ~5-10 pages.

\section{Datasets}

\subsection{ImageNet100}

ImageNet100 is a curated subset of the full ImageNet collection \cite{deng2009imagenet}, containing images from 100 object categories such as aeroplanes, cars, and chairs. For the purposes of this project, the collection is pre-organized into two subsets: InDSearch, containing approximately 95,000 images, and InDQuery, containing approximately 40,000 images, each covering the same 100 classes.

All images are resized to $224 \times 224$ pixels and passed through a WideResNet-50-2 neural network \cite{zagoruyko2016wide} to extract 2,048-dimensional feature vectors. Two network configurations are used: Conf1, which uses standard ImageNet-pretrained weights, and Conf2, which uses contrastively fine-tuned weights. The motivation for the two configurations is tied to the OoD detection methods described in Section~\ref{sec:ood_methods}. Feature extraction is computationally demanding, so all feature vectors were extracted once and stored on disk prior to the experiments to avoid repeated computation.

In the similarity search task, InDSearch serves as the database being searched, while InDQuery provides the query vectors. In the OoD detection task, InDSearch serves as the in-distribution reference set, InDQuery as in-distribution test queries, and Places365 as out-of-distribution test queries.

Places365 \cite{zhou2017places} is a collection of 36,500 scene images across 365 indoor and outdoor categories such as kitchens, streets, and offices. Since it depicts a fundamentally different type of visual content than ImageNet100, it serves as the out-of-distribution query set in the OoD experiments. The same preprocessing is applied: images are resized to $224 \times 224$ pixels and feature vectors are extracted using both Conf1 and Conf2.



\subsection{MovieLens 10M}



\section{Algorithms}

\subsection{Exact $k$-Nearest Neighbors}



\subsection{Locality Sensitive Hashing (LSH)}



\subsection{Hierarchical Navigable Small World (HNSW)}



\section{Out-of-Distribution Detection Methods}
\label{sec:ood_methods}

\subsection{Deep Nearest Neighbors}



\subsection{NN-Guide}



\section{Recommendation System}
\label{sec:rec_system}

\subsection{Item-Based $k$-Nearest Neighbors}



\section{Evaluation Metrics}



% ─────────────────────────────────────────────────────────────────────────────
\chapter{Results and Analyses}
% ─────────────────────────────────────────────────────────────────────────────
% All results must be explicitly referenced and discussed in text.
% Use \ref{} and \label{} for all figures and tables.

\section{Similarity Search Task}



\section{Out-of-Distribution Detection Task}



\section{Recommendation Task}



% ─────────────────────────────────────────────────────────────────────────────
\chapter{Discussion}
% ─────────────────────────────────────────────────────────────────────────────
% High-level interpretations and main takeaways across all results.



% ─────────────────────────────────────────────────────────────────────────────
\chapter{Conclusions and Perspectives}
% ─────────────────────────────────────────────────────────────────────────────
% Final remarks, limitations, and possible future extensions.
% Suggested length: ~1-2 pages.



% ─────────────────────────────────────────────────────────────────────────────
\bibliographystyle{plain}
\bibliography{references}

% ─────────────────────────────────────────────────────────────────────────────
\appendix
\chapter*{Appendix: Individual Report}
\addcontentsline{toc}{chapter}{Appendix: Individual Report}

\paragraph{Task Division} Not Applicable (individual project).

\paragraph{Name and Individual Contribution} Not Applicable.

\paragraph{Organizational Issues}
% Describe any organizational/communication issues encountered.

\paragraph{Went Better than Planned}
% Describe what you are particularly proud of.

\paragraph{Went Worse than Planned}
% What did not work well? What could be done differently?

\paragraph{Main Technical Takeaways}
% Main technical lessons learned and their relevance to your professional journey.

\paragraph{Other Takeaways}
% Non-technical lessons (organization, writing, reporting, etc.).

\paragraph{What If\ldots}
% If you had to do it again, what would you do differently?

\end{document}
